\documentclass[11pt]{article}
\newcommand{\DocShort}{Especificaciones t\'ecnicas}
\usepackage{fontspec}
\usepackage[spanish,es-noshorthands]{babel}
\decimalpoint
\usepackage{geometry}
\geometry{letterpaper, top=2.5cm, bottom=2.5cm, left=2.8cm, right=2.8cm}
\usepackage[expansion=false]{microtype}
\usepackage{parskip}
\setlength{\parskip}{6pt}
\usepackage{amsmath,amssymb}
\usepackage{booktabs}
\usepackage{array}
\usepackage{longtable}
\usepackage{caption}
\usepackage[dvipsnames,table]{xcolor}
\usepackage{enumitem}
\setlist[itemize]{topsep=2pt, itemsep=2pt, parsep=0pt}
\setlist[enumerate]{topsep=2pt, itemsep=2pt, parsep=0pt}
\usepackage{fancyhdr}
\usepackage[hidelinks,pdfencoding=auto]{hyperref}
\usepackage{titlesec}
\usepackage{tcolorbox}
\tcbuselibrary{breakable}
\usepackage{pdfpages}

\definecolor{darkblue}{HTML}{1B3A5C}
\definecolor{medblue}{HTML}{2E6DA4}
\definecolor{lightblue}{HTML}{D6E8F7}
\definecolor{teal}{HTML}{1A7A6E}
\definecolor{lightteal}{HTML}{D0EDEA}
\definecolor{lightgrey}{HTML}{F5F5F5}
\definecolor{midgrey}{HTML}{6F6F6F}
\definecolor{warnyellow}{HTML}{FFF3CD}

\titleformat{\section}
{\large\bfseries\color{darkblue}}
{\thesection}{0.75em}{}[{\color{medblue}\titlerule[0.8pt]}]
\titleformat{\subsection}
{\normalsize\bfseries\color{darkblue}}
{\thesubsection}{0.75em}{}
\titleformat{\subsubsection}
{\normalsize\bfseries\color{medblue}}
{\thesubsubsection}{0.75em}{}

\pagestyle{fancy}
\fancyhf{}
\fancyhead[L]{\small\color{midgrey}ICV 442 Acueductos y Alcantarillados}
\fancyhead[R]{\small\color{midgrey}\DocShort}
\fancyfoot[C]{\small\color{midgrey}\thepage}
\renewcommand{\headrulewidth}{0.4pt}
\renewcommand{\footrulewidth}{0pt}

\rowcolors{2}{lightgrey}{white}
\renewcommand{\arraystretch}{1.18}

\newtcolorbox{infobox}{
  breakable,
  colback=lightblue,
  colframe=medblue,
  boxrule=0.8pt,
  arc=3pt,
  left=8pt,
  right=8pt,
  top=7pt,
  bottom=7pt
}

\newtcolorbox{warnbox}{
  breakable,
  colback=warnyellow,
  colframe=orange!70!black,
  boxrule=0.8pt,
  arc=3pt,
  left=8pt,
  right=8pt,
  top=7pt,
  bottom=7pt
}

\newcommand{\doccover}[3]{
  \begin{center}
  {\small\color{midgrey}Pontificia Universidad Cat\'olica Madre y Maestra\\Escuela de Ingenier\'ia Civil y Ambiental}

  \vspace{1.0cm}
  {\color{medblue}\rule{\linewidth}{2pt}}
  \vspace{0.3cm}

  {\Large\bfseries\color{darkblue}ICV-442-T Acueductos y Alcantarillado}\\[0.35em]
  {\LARGE\bfseries\color{darkblue}#1}\\[0.25em]
  {\large\itshape #2}

  \vspace{0.3cm}
  {\color{medblue}\rule{\linewidth}{2pt}}
  \vspace{0.9cm}
  \end{center}

  \begin{center}
  \begin{tabular}{ll}
  \toprule
  \textbf{Asignatura} & ICV-442-T Acueductos y Alcantarillado \\
  \textbf{Profesor} & Ricardo Rafael Hern\'andez Moreira, PhD \\
  \textbf{Periodo acad\'emico} & Mayo--agosto de 2026 \\
  \textbf{Proyecto} & Dise\~no de sistemas de saneamiento para Las Charcas, Azua \\
  \textbf{Versi\'on} & #3 \\
  \bottomrule
  \end{tabular}
  \end{center}

  \vspace{0.4cm}
}


\begin{document}
\doccover{Especificaciones t\'ecnicas de dise\~no}{Criterios m\'inimos para el desarrollo del proyecto}{18 de mayo de 2026}

\begin{warnbox}
Estas especificaciones fijan criterios m\'inimos de dise\~no para el curso. No sustituyen la lectura de la normativa aplicable ni eximen al equipo de justificar los valores adoptados cuando existan alternativas aceptables.
\end{warnbox}

\section{Proyecci\'on poblacional y par\'ametros base}
\begin{center}
\begin{tabular}{p{0.28\linewidth} p{0.56\linewidth}}
\toprule
\textbf{Aspecto} & \textbf{Criterio} \\
\midrule
Fuente base & Datos oficiales de la ONE para Las Charcas, Azua \\
M\'etodo de proyecci\'on & Selecci\'on y justificaci\'on de un m\'etodo coherente con la serie hist\'orica disponible \\
Horizonte de dise\~no & 20 a\~nos para agua potable; 30 a\~nos para alcantarillado sanitario y pluvial \\
Soporte de c\'alculo & Tabla reproducible en hoja de c\'alculo con f\'ormulas visibles \\
Uso en el proyecto & La poblaci\'on proyectada debe alimentar demandas, aportes y estimaciones hidrol\'ogicas \\
\bottomrule
\end{tabular}
\end{center}

\section{Agua potable}
\subsection{Criterios generales}
\begin{center}
\begin{tabular}{p{0.27\linewidth} p{0.57\linewidth}}
\toprule
\textbf{Tema} & \textbf{Criterio} \\
\midrule
Fuente de abastecimiento & Debe justificarse la alternativa adoptada a partir de disponibilidad, calidad y factibilidad \\
P\'erdidas de carga & El proyecto usar\'a Darcy-Weisbach como m\'etodo base para conducciones y red \\
Presi\'on en servicio & Mantener rangos compatibles con la normativa aplicable y el buen funcionamiento de la red \\
Velocidades & Verificar que no sean ni insuficientes ni excesivas para las tuber\'ias propuestas \\
Accesorios principales & Incluir, cuando corresponda, v\'alvulas, purgas, ventosas, hidrantes y elementos de control \\
\bottomrule
\end{tabular}
\end{center}

\subsection{Modelaci\'on en EPANET}
\begin{itemize}
\item Asignar demandas de forma coherente con la distribuci\'on espacial de la poblaci\'on.
\item Cargar el patr\'on horario suministrado en el anexo correspondiente.
\item Reportar presiones, velocidades y nodos cr\'iticos.
\item Incorporar, cuando proceda, un escenario de verificaci\'on adicional para demanda extraordinaria o incendio, seg\'un el alcance definido por el equipo y la normativa consultada.
\end{itemize}

\section{Alcantarillado sanitario}
\subsection{Criterios generales}
\begin{center}
\begin{tabular}{p{0.27\linewidth} p{0.57\linewidth}}
\toprule
\textbf{Tema} & \textbf{Criterio} \\
\midrule
Caudales de dise\~no & Deben construirse a partir de poblaci\'on, contribuci\'on per c\'apita, factores de pico y aportes complementarios justificados \\
Infiltraci\'on y conexiones erradas & Deben considerarse solo en la medida en que la norma o el criterio adoptado lo exijan; el valor asumido debe explicarse \\
Velocidades y pendientes & Verificar autolimpieza, capacidad y profundidad de zanja dentro de rangos t\'ecnicamente razonables \\
Materiales & Selecci\'on abierta, siempre que se justifique por desempe\~no, disponibilidad y adecuaci\'on al contexto \\
Bombeo & Solo cuando la topograf\'ia o la configuraci\'on del sistema lo haga necesario \\
\bottomrule
\end{tabular}
\end{center}

\subsection{Modelaci\'on en SWMM para sanitario}
\begin{itemize}
\item Cargar el patr\'on horario de descarga sanitaria suministrado en el anexo.
\item Introducir los aportes de tiempo seco de manera clara y consistente con la memoria.
\item No es necesario introducir m\'etodos de infiltraci\'on hidrol\'ogica de lluvia para la estimaci\'on del caudal sanitario de redes separativas.
\item Reportar niveles, velocidades y tramos cr\'iticos.
\end{itemize}

\section{Drenaje pluvial}
\subsection{Hidrolog\'ia de dise\~no}
\begin{center}
\begin{tabular}{p{0.27\linewidth} p{0.57\linewidth}}
\toprule
\textbf{Tema} & \textbf{Criterio} \\
\midrule
Periodo de retorno & Seleccionar seg\'un jerarqu\'ia vial, uso del suelo y consecuencia de falla \\
M\'etodo hidrol\'ogico & Racional para cuencas peque\~nas o m\'etodo alternativo justificable con la informaci\'on disponible \\
Tiempo de concentraci\'on & Elegir y justificar el m\'etodo empleado \\
Curvas IDF & Documentar la procedencia y el procedimiento de interpolaci\'on o ajuste \\
Subcuencas & Delimitar con apoyo del MDT y del reconocimiento del terreno \\
\bottomrule
\end{tabular}
\end{center}

\subsection{Dise\~no hidr\'aulico}
\begin{itemize}
\item Verificar capacidad, velocidades, pendientes y control de descargas.
\item Definir con claridad puntos de captaci\'on y descarga.
\item Considerar separaci\'on adecuada respecto del sistema sanitario.
\item Incluir estructuras de disipaci\'on o protecci\'on cuando las condiciones de descarga lo exijan.
\end{itemize}

\subsection{Modelaci\'on en SWMM para pluvial}
\begin{itemize}
\item Construir subcuencas y nodos de forma consistente con el trazado adoptado.
\item Usar un evento de dise\~no expl\'icito y documentado.
\item Reportar caudales pico, niveles y puntos con comportamiento cr\'itico.
\end{itemize}

\section{Contenido m\'inimo de la memoria t\'ecnica}
\begin{itemize}
\item Tabla maestra de par\'ametros con referencia normativa o bibliogr\'afica.
\item Desarrollo de c\'alculos principales.
\item Resultados de modelaci\'on con tablas y figuras legibles.
\item Justificaci\'on de materiales y criterios constructivos.
\item Comentario final sobre factibilidad, limitaciones y mantenimiento.
\end{itemize}

\section{Criterios de aceptaci\'on}
\begin{enumerate}
\item Coherencia entre la informaci\'on disponible, los supuestos y el dise\~no propuesto.
\item Justificaci\'on normativa suficiente en par\'ametros clave.
\item Modelos que permitan verificar el desempe\~no hidr\'aulico del sistema.
\item Documentaci\'on clara, legible y trazable.
\item Compatibilidad entre agua potable, sanitario y pluvial.
\end{enumerate}

\end{document}
